 \documentclass[12pt]{article}
\usepackage[a4paper,margin=2.2cm]{geometry}
\usepackage{amsmath,amssymb,amsthm,mathtools}
\usepackage{hyperref}
\usepackage{cleveref}

% 中文（xelatex）
\usepackage{fontspec}
\usepackage{xeCJK}

\usepackage{fontspec}
\usepackage{xeCJK}
\setCJKmainfont{Noto Serif CJK SC}

\title{Linear Algebra-Ega2}
\date{}
\begin{document}
\maketitle


\section{破裂日 (Defective Matrix)}
\begin{flushleft}
    \[ A = \begin{bmatrix}
            2 & 1 \\
            0 & 2
        \end{bmatrix} \]
    尝试对其进行对角化：

    \vspace{0.5em}
    \noindent\textbf{Step 1: Solve $\lambda$ (求特征值与代数重数)} \\
    \[
        \det(A - \lambda I) = 0
    \]
    \[
        (2-\lambda)^{2}  = 0
    \]
    解得特征值为重根 $\lambda = 2$。\\
    因为因子的指数为 2，所以：\textbf{代数重数 = 2}。

    \vspace{0.5em}
    \noindent\textbf{Step 2: Solve egv (求特征向量与几何重数)} \\
    将 $\lambda = 2$ 代入 $A - \lambda I$：
    \[
        \begin{bmatrix}
            0 & 1 \\
            0 & 0
        \end{bmatrix}
    \]
    显然，该矩阵的 Rank (秩) = 1。\\
    根据定理：\textbf{几何重数} = 维度 $n$ - Rank = $2 - 1 = \textbf{1}$。

    求解齐次方程组 $(A - 2I)x = 0$：
    \[ \begin{cases}
            0x  +  y = 0 \\
            0x  +  0y = 0
        \end{cases} \]
    解得：$y = 0$，$x$ 可以为任何非零数。\\
    得到\textbf{唯一}的基础解系（特征向量）：
    $\vec{e}_1 =  \begin{bmatrix} 1 \\ 0 \end{bmatrix}$

    \vspace{0.5em}
    \noindent\textbf{Step 3: 为什么 $g_{i} < a_{i}$ 不能对角化？} \\
    对角化的核心公式是：$A = P \cdot D \cdot P^{-1}$。\\
    其中矩阵 $P$ 必须是由 $n$ 个特征向量拼成的可逆矩阵：$P = \begin{bmatrix} \vec{e}_1 & \cdots & \vec{e}_n \end{bmatrix}$。

    要让 $P$ 可逆（即能算出 $P^{-1}$），$P$ 必须满足：
    \begin{enumerate}
        \item 内部的列向量必须\textbf{线性无关}（例如不能存在 $\vec{e}_i = k \cdot \vec{e}_j$ 的共线情况）。
        \item 矩阵必须满秩，即 \textbf{$\det(P) \neq 0$}（不能出现降维塌陷）。
    \end{enumerate}


    \vspace{0.5em}
    \noindent\textbf{Step 4: 到底什么是对角化？它和 $P$ 以及 $P^{-1}$ 的几何关系} \\
    对角化就是找到矩阵 $A$ 的“极简形态” $D$（由特征值组成的对角矩阵）。

    \textbf{代数推导：} \\
    如果 $P$ 可逆，对角化公式为：$A = P \cdot D \cdot P^{-1}$。\\
    左右同乘 $P$：$AP = P \cdot D \cdot P^{-1} P \implies AP = PD$ \\
    左乘 $P^{-1}$：$P^{-1}AP =  P^{-1}PD \implies D = P^{-1}AP$

    \textbf{几何直觉（翻译器理论）：} \\
    我们把一个普通向量 $\vec{x}$ 放到公式最右边来观察变换过程：$\textbf{A}\vec{x} = \textbf{P} \cdot \textbf{D} \cdot \textbf{P}^{-1} \vec{x}$。\\
    矩阵乘法是\textbf{从右向左}执行的，这完美对应了三个步骤：
    \begin{enumerate}
        \item \textbf{$P^{-1}$ (入境翻译器)：} 把向量 $\vec{x}$ 从“标准基环境”翻译并拉拽到“$P$ 特征向量基环境”中。
        \item \textbf{$D$ (极简处理核心)：} 在特征向量的专属世界里，变换变得极其简单！不需要复杂的旋转扭曲，只需要沿着各个轴做简单的缩放（乘以特征值）。这就是到达目的地的捷径。
        \item \textbf{$P$ (出境翻译器)：} 处理完毕后，用 $P$ 把结果从“特征环境”重新翻译回我们熟悉的“标准基环境”。
    \end{enumerate}

    从几何视角分析，\textbf{$A$ 是基于标准基坐标系的、极其混乱复杂的转换；而 $D$ 是基于特征向量基坐标系的、最为简单的转换（只有单纯的拉伸/压缩）。} \\
    但无论走哪条路（直接用 $A$，或者走 $P \rightarrow D \rightarrow P^{-1}$ 的捷径），它们最终都会到达\textbf{同一个绝对的向量地址}。对角化，就是为了寻找这条“降维打击”般的捷径！

    % \noindent\textbf{Step 4: 到底什么称为对角化 ? 它和 P和 $P^{-1}$ 的关系}: \\
    % 对角化 = D, D 是对角矩阵由特征值组成 \\
    % 如果 P 可逆:
    % $A = P \cdot D \cdot P^{-1}$。\\
    % $AP = P \cdot D \cdot P^{-1} P $ \\
    % $AP = PD$ \\
    % $P^{-1}AP =  P^{-1}PD$ \\
    % $P^{-1}AP = D$ \\
    % $D = P^{-1}AP$
    % 我们思考这个公式 $D = P^{-1}AP$ : \\
    % P 为 A 矩阵的翻译器(把向量翻译回 A 标准基环境) \\
    % $P^{-1}$ 是 P 矩阵的翻译器 (把向量翻译回 P 特征向量基环境) \\
    % D 就是 在 P 特征向量基环境 达到目的地的路径  因为 D 是 对角矩阵 所以它足够简单 \\
    % 所以我们从几何视角分析来说 A 是基于标准基坐标系的混乱的转换  而D 是基于特征向量基坐标系最为简单的转换 \\
    % 它们都可以到达同一个绝对向量地址

    \noindent\textbf{Step 5: 这个矩阵的几何行为是什么}: \\
    从A矩阵两个列向量初看 :
    应该是先拉伸之后向右上提 ?
    因为 $\vec{e1} = \begin{bmatrix}
            2 \\
            0
        \end{bmatrix} = 2 \begin{bmatrix}
            1 \\
            0
        \end{bmatrix}$ 是非常标准的拉伸动作
    但  $\vec{e2}.x = 1$ 而非0 则证明路径整体上移?

\end{flushleft}



\section{代数重数 与 几何重数}
\begin{flushleft}
    已知矩阵 $A = \begin{bmatrix} 2 & 0 \\ 0 & 2 \end{bmatrix}$，判断其是否可对角化。

    \vspace{0.5em}
    \noindent\textbf{第一步：求特征值与代数重数} \\
    构建特征方程：$\det(A - \lambda I) = 0$
    \[
        \begin{vmatrix}
            2 - \lambda & 0           \\
            0           & 2 - \lambda
        \end{vmatrix} = 0
    \]
    化简得多项式：$(2 - \lambda)^2 = 0$ \\
    解得特征值为 $\lambda = 2$。
    \begin{itemize}
        \item \textbf{代数重数：} 看特征多项式中 $(\lambda - 2)$ 这个因子的\textbf{最高次幂}。因为这里是平方（2次方），所以 $\lambda = 2$ 的\textbf{代数重数 = 2}。
    \end{itemize}

    \vspace{0.5em}
    \noindent\textbf{第二步：求特征向量与几何重数} \\
    将 $\lambda = 2$ 代入，构建齐次线性方程组 $(A - \lambda I)x = 0$：
    \[
        \begin{bmatrix} 0 & 0 \\ 0 & 0 \end{bmatrix}
        \begin{bmatrix} x_1 \\ x_2 \end{bmatrix} =
        \begin{bmatrix} 0 \\ 0 \end{bmatrix}
    \]
    展开即为：
    \[
        \begin{cases}
            0x_1 + 0x_2 = 0 \\
            0x_1 + 0x_2 = 0
        \end{cases}
    \]

    \begin{itemize}
        \item \textbf{计算几何重数（重点）：} \\
              此时系数矩阵 $(A - 2I)$ 是一个全零矩阵，其\textbf{秩 $r = 0$}。\\
              矩阵的阶数 $n = 2$（未知数个数）。\\
              根据定理：特征空间的维数（几何重数） $= n - r$ \\
              所以，\textbf{几何重数 = $2 - 0 = 2$}。
    \end{itemize}

    \vspace{0.5em}
    \noindent\textbf{第三步：求基础解系} \\
    因为方程组的秩为 0，说明 $x_1$ 和 $x_2$ 都是自由变量（可以取任意不全为 0 的数）。\\
    我们可以取两组线性无关的向量作为基础解系：
    \[
        p_1 = \begin{bmatrix} 1 \\ 0 \end{bmatrix}, \quad p_2 = \begin{bmatrix} 0 \\ 1 \end{bmatrix}
    \]
    这说明特征空间是二维的。

    \vspace{0.5em}
    \noindent\textbf{结论：} \\
    对于特征值 $\lambda=2$，其 \textbf{几何重数(2) = 代数重数(2)}。\\
    因此，矩阵 $A$ 可以找到 2 个线性无关的特征向量，\textbf{该矩阵可对角化}（实际上它本身就是对角阵）。



    问一个一直困扰我很久的模糊问题 为什么矩阵乘法的运算规则是行*列
    $A = \begin{bmatrix}
            a & b \\
            c & d
        \end{bmatrix}$ \\
    我在理解矩阵时 会以列向量的模式理解 也就是矩阵 A 是有 $e1 = \begin{bmatrix}
            a \\
            c
        \end{bmatrix},e2 = \begin{bmatrix}
            b \\
            d
        \end{bmatrix}$ 组成 是基底向量 我们知道 标准矩阵 $I = \begin{bmatrix}
            1 & 0 \\
            0 & 1 \\
        \end{bmatrix}$ A 和 I 的变化关系是 $ I.e1 -> A.e1  I.e2 -> A.e2$ 对于一个向量 $\vec{v}$ \\
    经过 A 矩阵变化之后 -> $\vec{v_{n}} = v.x \cdot A.e1  + v.y \cdot A.e2 $ 一种向量加法的形态
    但这个代数形态并不是行*列呀 ?


    % 你的意思是我的思路和理解是正确的 只是少了一次代数推理 ?
    % 你的回复有个点很有意思 你说矩阵诞生的目的是为了解决 解线性方程组
    % 问题 1: 为什么线性方程组能使用矩阵这形态表达? 逻辑是什么? 
    % 问题 2: 并不是为了做空间几何变换，而是为了解线性方程组 ? 为什么后来矩阵可以做空间几何变换呢 ? 是什么特性或发现 让矩阵成为现在的矩阵 而非只是解线性方程组的另一种形态呢 ?
\end{flushleft}

\section{矩阵}
\begin{flushleft}

    \vspace{0.5em}
    \noindent\textbf{1: 矩阵的诞生与矩阵的代数意义?} \\
    为什么课题是矩阵与向量 矩阵在前向量在后? 因为矩阵的诞生是为了解决代数问题: \\
    存在方程组 a :\\
    \[
        \begin{cases}
            4x + 3y = 8 \\
            4x - y  = 5
        \end{cases}
    \]
    a 存在一个重要的代数性质 "线性" 也就是满足可加性、齐次性 \\
    假如我们把方程组 a 的 左边 抽象成函数 f(x, y) : \\
    可加性满足: \\
    $f(u+v) = f(u) + f(v)$ \\
    齐次性满足:
    $f(k \cdot u) = k \cdot f(u)$ \\

    a 方程可以变化成一个矩阵模式(这里只有代数意义) \\
    $ A =  \begin{bmatrix}
            4 & 3  \\
            4 & -1 \\
        \end{bmatrix} \begin{bmatrix}
            x \\
            y
        \end{bmatrix} = \begin{bmatrix}
            8 \\
            5
        \end{bmatrix}$ \\
    A 可以通过矩阵乘法规则还原成 a

    \vspace{0.5em}
    \noindent\textbf{2: 矩阵模式对于方程组的优势(数学家为什么要发明矩阵这个概念去表达代数方程组) ? } \\

    A 是可以还原为 a 的 利用矩阵乘法规则 \\
    那么我们仔细分析 A 的形态:
    $ A\cdot x  = y$ \\
    X 完全为代数未知数, x 和 y 是完全的已知数 \\
    1: 把未知和已知解耦 这里是利用了齐次性 $f(k \cdot x) = k \cdot f(x)$ \\ k 是已知数 f(x) 是未知数 ? \\
    2: 如果 Y 的系数项 非常大 那么整个代数运算也会是 $X \cdot Y  = Z$  这种形态 X 依然简单 Y 包含了全部的线性运算规则 隔离成了两个整体 \\

    \vspace{0.5em}
    \noindent\textbf{3: 矩阵和线性的关系(可加性、齐次性)} \\
    A 是可以还原为 a 所以 矩阵形态 是天然满足线性的(或者说 A 是容器, 矩阵乘法规则才是线性的表达)\\
    $A(u+v)  = Au + Av$ 和 $A(ku)  = k \cdot Au$


    \vspace{0.5em}
    \noindent\textbf{4: 矩阵和几何空间的关系,为什么矩阵变成了可以表达空间几何变换？是什么特性造成的？} \\
    我们由上已知道 矩阵形态是天然满足线性的 \\
    所以我们关键的要分析 可加性 和 齐次性 在几何上是什么形态表达 ? \\
    齐次性的几何表现:  \\
    $f(k \cdot v) = k \cdot f(v)$ \\
    齐次性是使用代数乘法的运算 如果 v 是一个有方向的数学对象 那么 $k \cdot v$ = $v + v + \cdots$ \\
    我们会发现 它会沿这个方向 (如果 k 是负数 就是反方向) 做以 v 为单位的伸缩动作 这个动作就是直线伸缩动作 \\
    v 有起点吗 ? \\
    我们把 k 设置为 0 ,这就是代数上的起点  \\
    $f(0 \cdot v) = 0 \cdot f(v) = 0$ \\
    在几何上这个起点就是零空间 也就是原点 \\
    我们需要理解上面这个代数式 f 是线性变化  $0 \cdot v$ 是处于原点的 v \\
    $= 0 \cdot f(v)$ 如果满足齐次性 这个等号满足 那么结果必然等于 0 \\
    所以我们能得到 齐次性的几何表现是 过原点的 直线伸缩动作 ?
    只要一个变换满足齐次性，它就绝对不能移动原点，原点必须是变换的不动点。反过来，所有会移动原点的变换（比如平移），必然不满足齐次性，也就不是线性变换
    可加性的几何表现:  \\
    $f(u+v) = f(u) + f(v)$ \\
    可加性是使用代数加法而运算 加法有一个非常意思的特性 -> 支持对同一个结构的不同数据的多个数学对象进行组合平移(例如 自然数相加 3 + 4 = 7 他们是同一个结构但值不一样 \\
    例如 $\vec{v} + \vec{x} = \vec{vx} $) \\
    而乘法是复制平移是单方向的 但组合平移是多方向的(或者说是不限制方向的) \\
    $f(u+v) = f(u) + f(v) = f(v) + f(u)$ \\
    这里我们拿路线为例子 f(u) 和 f(v) 都是一个拥有路线(有方向)的目的地 先走 f(u) 之后再走 f(v) 到达的最终目的地  = 先走 f(v) 再 走 f(u) \\
    是天然存在对称性的(加法也存在对称性) 那么什么几何形状存在两步结构天然对称性的? 平行四边形 \\
    所以 可加性是 任意方向组合平行四边形? \\

    所以当矩阵形态(重点是矩阵运算规则)是满足线性的 那么矩阵天然适配为线性几何变化, 它只是被发现 \\

\end{flushleft}


\section{逆向问题}
\begin{flushleft}
    已知 $D = \begin{bmatrix}
            3 & 0  \\
            0 & -1
        \end{bmatrix}, P = \begin{bmatrix}
            1 & 1  \\
            1 & -1
        \end{bmatrix}$

    \vspace{0.5em}
    \noindent\textbf{1: 求 A} \\
    $A  = P\cdot D \cdot P^{-1}$ \\
    $det(P) = -1 - 1 = -2$ \\
    $P^{-1} = \frac{1}{ad-bc}\begin{bmatrix}
            d  & -b \\
            -c & a
        \end{bmatrix}  = \begin{bmatrix}
            \frac{1}{2} & \frac{1}{2}  \\
            \frac{1}{2} & -\frac{1}{2}
        \end{bmatrix}$ \\
    $A = \begin{bmatrix}
            1 & 1  \\
            1 & -1
        \end{bmatrix}  \begin{bmatrix}
            3 & 0  \\
            0 & -1
        \end{bmatrix} \begin{bmatrix}
            \frac{1}{2} & \frac{1}{2}  \\
            \frac{1}{2} & -\frac{1}{2}
        \end{bmatrix} =  \begin{bmatrix}
            1 & 2 \\
            2 & 1
        \end{bmatrix}$

    \vspace{0.5em}
    \noindent\textbf{2: A 的特征值是否真的是 3 和 -1} \\
    迹 + 特征多项式 / 韦达定理
    \begin{align}
        \lambda_{1} + \lambda_{2} = 2 \;(1 + 1) \\
        \lambda_{1} \cdot \lambda_{2} = -3 \; (1 \cdot  1  - 2 \cdot 2 )
    \end{align}
    带入 3, -1 成立

    \vspace{0.5em}
    \noindent\textbf{3: 思考题：你是从"特征基世界"出发，构造出标准基下的矩阵。用你的翻译器类比说清楚这个反向过程} \\
    $A\vec{v}  = P\cdot D \cdot P^{-1} \vec{v}$ \\
    我们已知 P 和 D 是矩阵 特征向量矩阵 和特征值对角矩阵 \\
    所以我们依然重点要关注 P 到底是什么: \\
    P 最重要的一个特点是 它使用基于基础基底(e1,e2)的两个特征向量 同样撑起了整个以基础基底(e1,e2)的平面 \\
    $P = \begin{bmatrix}
            \vec{g1} | \; \vec{g2}
        \end{bmatrix}$

    我们以一个最终向量$\vec{X} = \begin{bmatrix}
            p1 \\
            p2
        \end{bmatrix}$

    那么存在下面的代数关系(也是几何关系) : \\
    $\vec{v} = \begin{bmatrix}
            a1 \\
            a2
        \end{bmatrix},\vec{v1} = \begin{bmatrix}
            c1 \\
            c2
        \end{bmatrix} $ \\
    $= a1 \cdot e1 + a2  \cdot e2  = c1 \cdot g1  + c2 \cdot g2  = \vec{X}$ \\
    a1,a2,c1,c2 为常量 代表步伐 \\

    g1, g2 可以进一步展开为  \\
    $g1 = g1_{x}e1 + g1_{y}e2$ \\
    $g2 = g2_{x}e1 + g2_{y}e2$ \\
    代数转换一下 :\\
    $= c1 \cdot (g1_{x}e1 + g1_{y}e2)  + c2 \cdot (g2_{x}e1 + g2_{y}e2)$ \\

    $= \begin{bmatrix}
            \vec{e1} | \; \vec{e2}
        \end{bmatrix} c1 \cdot \begin{bmatrix}
            g1_{x} \\
            g1_{y}
        \end{bmatrix}  + \begin{bmatrix}
            \vec{e1} | \; \vec{e2}
        \end{bmatrix}  c2 \cdot \begin{bmatrix}
            g2_{x} \\
            g2_{y}
        \end{bmatrix} $ \\

    $ = \begin{bmatrix}
            \vec{e1} | \; \vec{e2}
        \end{bmatrix}(c1 \cdot \begin{bmatrix}
            g1_{x} \\
            g1_{y}
        \end{bmatrix} + c2 \cdot \begin{bmatrix}
            g2_{x} \\
            g2_{y}
        \end{bmatrix}) $ \\


    $\begin{bmatrix}
            \vec{e1} | \; \vec{e2}
        \end{bmatrix}\begin{bmatrix}
            a1 \\
            a2
        \end{bmatrix} = \begin{bmatrix}
            \vec{e1} | \; \vec{e2}
        \end{bmatrix} (c1 \cdot \begin{bmatrix}
            g1_{x} \\
            g1_{y}
        \end{bmatrix} + c2 \cdot \begin{bmatrix}
            g2_{x} \\
            g2_{y}
        \end{bmatrix})$ \\


    $\begin{bmatrix}
            a1 \\
            a2
        \end{bmatrix} = c1 \cdot \begin{bmatrix}
            g1_{x} \\
            g1_{y}
        \end{bmatrix} + c2 \cdot \begin{bmatrix}
            g2_{x} \\
            g2_{y}
        \end{bmatrix} $ \\
    $\begin{bmatrix}
            a1 \\
            a2
        \end{bmatrix} = P \begin{bmatrix}
            c1 \\
            c2
        \end{bmatrix}$

    P 的列向量（特征向量）和标准基张成的是同一个$R^{2}$空间, 把一个向量翻译回这个空间 \\


    $P^{-1}$ 到底是什么 ? \\
    我们已知 P 是一个翻译器  那么 $P^{-1}$ 依然是一个翻译器 但我们要使用代数证明 $P^{-1}$ 翻译到什么环境?? \\
    同时乘法 $P^{-1}$

    $P^{-1}\begin{bmatrix}
            a1 \\
            a2
        \end{bmatrix} =P^{-1} P \begin{bmatrix}
            c1 \\
            c2
        \end{bmatrix}$ \\
    $= P^{-1}\begin{bmatrix}
            a1 \\
            a2
        \end{bmatrix} =  \begin{bmatrix}
            c1 \\
            c2
        \end{bmatrix} $ \\
    $ P^{-1}\begin{bmatrix}
            a1 \\
            a2
        \end{bmatrix} -> \begin{bmatrix}
            c1 \\
            c2
        \end{bmatrix}$  这个代数式非常清楚的表达了 $P^{-1}$ 的身份 把一个向量翻译到特征向量基底空间  \\


    那么这个逆向工程就非常清晰了 : 我们拥有 P 和 D 我们可以找到 $P^{-1}$

    $DP^{-1}$ 是把向量翻译到特征向量基底空间并到达目的地,   P 把这个目的地翻译回 A(P 和 A 在一个张成空间)

\end{flushleft}

\section{矩阵的迹}
\begin{flushleft}
    定义 : 矩阵$A \in \mathbb{R}^{n \times n}$ \\
    $tr(A) = \sum_{i = 1}^{n} A_{ii}$
    \begin{itemize}
        \item 1: 到底什么是迹, 它是为了解决什么问题 ? \\
              在标准基下(e1, e2): \\
              $A = \begin{bmatrix}
                      a & b \\
                      c & d
                  \end{bmatrix}$\\
              我们要分析一个问题: a, b, c, d 这些常数项的变化会给整个矩阵带来什么变化 ?  \\
              $A\vec{v}  = \begin{bmatrix}
                      a & b \\
                      c & d
                  \end{bmatrix}\begin{bmatrix}
                      x \\
                      y
                  \end{bmatrix} = \begin{bmatrix}
                      ax + by \\
                      cx + dy
                  \end{bmatrix}$ \\
              我们设 A = $\begin{bmatrix}
                      1 & 0 \\
                      0 & 1
                  \end{bmatrix}$ \\
              $ \begin{bmatrix}
                      ax + by \\
                      cx + dy
                  \end{bmatrix} = \begin{bmatrix}
                      x \\
                      y
                  \end{bmatrix}$ \\
              我们改变 a.b != 0 \\
              $ = \begin{bmatrix}
                      x + by \\
                      y
                  \end{bmatrix}$
              所以我们发现 新的 y 依然是 y  新的 x 是 = x+by :
              例子: $\begin{bmatrix}
                      1 \\
                      1
                  \end{bmatrix}(v_{1}) ->  \begin{bmatrix}
                      b \cdot (1) \\
                      1
                  \end{bmatrix}(v_{2})$  $v_{2}$ 对比 $\vec{1}$ 发生了水平滑移 所以 b 控制了 水平剪切 程度 \\
              c:  \\
              $ = \begin{bmatrix}
                      x \\
                      cx + y
                  \end{bmatrix}$  发生了垂直剪切  \\
              a:
              $ = \begin{bmatrix}
                      ax \\
                      y
                  \end{bmatrix}$ 拉伸了 x 轴长 是水平拉伸 \\
              d:
              $ = \begin{bmatrix}
                      x \\
                      by
                  \end{bmatrix}$ 拉伸了 y 轴长 是垂直拉伸 \\
              我们现在知道 a, b 是倍数关系 : \\
              $tr(A) = a + d $ \\
              我们就能清晰的明白:  a + d  或 $\sum_{i = 1}^{n} A_{ii}$ 是剔除 剪切变化之后的 总伸缩率(这就是一种解耦的表现 -> 改变面积变化和改变形状变化)
              迹到底解决了什么问题: 对矩阵变化之后的伸缩率的一种直观表现?
        \item 2 迹与基的关系 : \\
              2.1 在回答迹与基的关系之前, 我们来弄明白什么叫换基和换基的目的 ? \\
              存在一个绝对的地址 可以使用绝对向量 X 表达 \\
              A 矩阵是基于标准基到达 X 的线性变换 \\
              B 矩阵是基于 P(基底)到达 X 的 线性变换 \\
              这是同一种抽象的线性变化, 在不同基下的矩阵表示 \\
              子问题2.1: B 的代数表达是什么(基于 A 和 P), 为什么这么表达 ? \\
              2.1.1 : $P$ 是什么, $Pv$ 是什么  ? \\
              A 是基于标准基底 E 的线性变换 $\begin{bmatrix}
                      \vec{e1} | \; \vec{e2}
                  \end{bmatrix}$ \\
              向量 $\vec{m} = \begin{bmatrix}
                      x \\
                      y
                  \end{bmatrix}$ 存在于标准基环境  \\
              向量 $\vec{v} = \begin{bmatrix}
                      c1 \\
                      c2
                  \end{bmatrix}$ \\
              因为 P 的身份是基底  $\begin{bmatrix}
                      \vec{p1} | \; \vec{p2}
                  \end{bmatrix}$\\
              但 P 还有另一种非常重要的特性: 它本身是基于 E 的一个抽象线性变换
              也就是说 如果 P 是存在的 那么它也能张成整个空间 \\
              那么对 X(绝对向量来说) \\
              $X = xe1 + ye2 =  c1p1 + c2p2$ \\
              $m = Pv $ \\
              那么 Pv 就是 m 是同一个对象, 不同的视角, P 是翻译回 E 环境的翻译器 \\
              2.1.2: B 的代数表达 ;
              $Am = APv$ 是成立的(基于 2.1.1) \\
              它们是 E 环境下完全相同的变换 \\
              已知 B 是 P 基底的线性变换  \\
              我们需要把 \\
              $P^{-1}Am = P^{-1}APv$ \\
              $ m = Pv$ \\
              上面的代数证明来 $P^{-1}$ 就是翻译回 P 环境的 翻译器！
              那么 $B = P^{-1}AP $

              有 2.1.1和 2.1.2 我们得知 B 和 A 是同一种线性变换 T 的 两种表达
              对 $tr(T) = tr(A) = tr(B)$ 迹在不同基底的不变性
        \item 3 迹与连续变化 XY : \\
              矩阵乘法并不支持交换率 $XYv != YXv$  \\
              我们也知道迹是线性的总伸缩率(伸缩率会改变面积, 但剪切变化不会改变面积, 只会改变形状(item1)) \\
              那么 $tr(XY)$ 的一种重要的代数性质是 先伸缩一个系数然后再伸缩一个系数 是一个连续动作 \\
              类似代数式 $s \cdot x \cdot y = s \cdot y  \cdot x$ (s 是面积), 这是一种逻辑推导 \\
              代数推导： \\
              $X,Y \in \mathbb{R}^{n \times n}$ \\
              X 的 i 行, j 列的数 记 $X_{ij}$ \\
              Y 的 j 行, k 列的数 记 $Y_{jk}$ \\
              step1: $tr(XY):$
              \begin{align}
                  (XY)ik = \sum_{j = 1}^{n} X_{ij}Y_{jk} \\
                  (XY)ii = \sum_{j = 1}^{n} X_{ij}Y_{ji} \\
                  tr(XY) = \sum_{i = 1}^{n}(\sum_{j = 1}^{n} X_{ij}Y_{ji})
              \end{align}
              Y 的 j 行, k 列的数 记 $Y_{jk}$ \\
              X 的 i 行, j 列的数 记 $X_{ij}$ \\
              step2: $tr(YX):$
              \begin{align}
                  (YX)ik = \sum_{i = 1}^{n} Y_{ij}X_{jk} \\
                  (YX)jj = \sum_{i = 1}^{n} Y_{ji}X_{ij} \\
                  tr(YX) = \sum_{j = 1}^{n}(\sum_{i = 1}^{n} Y_{ji}X_{ij})
              \end{align}
              $\sum_{j = 1}^{n}(\sum_{i = 1}^{n} Y_{ji}X_{ij})$ 是基于常数的加法与乘法的混合运算 是支持交换率和结合律的 \\
              $ = \sum_{j = 1}^{n}(\sum_{i = 1}^{n} X_{ij} Y_{ji})$ \\
              $ = \sum_{i = 1}^{n}(\sum_{j = 1}^{n} X_{ij} Y_{ji}) = tr(XY)$ \\
        \item 4 为什么 $tr(A) = \lambda_{1} + \lambda_{2}$ ? \\
              A 可以对角化
              $tr(A) = tr(PDP^{-1})$ \\
              $tr(A) = tr(XY) = tr(YX) = tr(P^{-1}PD) = tr(D) = \lambda_{1} + \lambda_{2}+ ...$ \\
              但我们也可以使用另一种视角 也就是 迹在不同基底的不变性 \\
              $tr(A) = tr(B) = tr(P^{-1}AP) =  tr(P^{-1}PDP^{-1}P) = tr(D) $
        \item 5 思考 为什么 $det(A) =  \lambda_{1} \cdot \lambda_{2}$ ? \\
              $det(A) = det(PDP^{-1})$\\
              $ = det(XY) = det(X)det(Y)\;\; (item3 \; 同样的代数证明逻辑)$ \\
              $ = det(P)det(DP^{-1}) = det(P)det(D)det(P^{-1}) = (使用交换律) det(P)det(P^{-1})det(D) = det(D)$

    \end{itemize}

\end{flushleft}

\end{document}
